% GENERATED FILE. Edit canonical lessons, then run npm run generate:books.
% canonical-source-hash: fnv1a64:e46a596dcffd42d5
% canonical-lessons: ZH-W23-de, ZH-C23-de, ZH-C23-de-mingci, ZH-C23-huati, ZH-C23-practice

\chapter{De — The Hinge}
\label{ch:zh-de}
\hlchaptermodality{23}

\begin{chapteropening}
\textbf{By the end of this chapter:} \emph{I can say whose something is, put a whole description in front of the noun it describes, and open a sentence with its topic.}

One character, eight strokes and no tone, doing three jobs — and the third of them, the topic slot, costs no character at all because it is an order rather than a word.
\end{chapteropening}

\section[de]{\zh{的} — eight strokes for the commonest character there is}

\label{lesson:ZH-W23-de}



\begin{quote}
A new chapter. Six retrievals, then the pen.

\begin{itemize}
\raggedright
  \item \emph{Recall:} say \emph{Zhōngguó rén — shū — kàn} — \textbf{R1}, one lesson back, before it decays
  \item \emph{Recall:} say \emph{kàn} — \textbf{R2}, the first real retrieval, five to fifteen lessons back
  \item \emph{Recall:} say \emph{kànshū} — \textbf{R3}, durable, twenty to sixty lessons back
  \item \emph{Recall:} say \emph{kàn — kànjiàn — hǎokàn — shū — kànshū} — \textbf{R3}, durable again, a second word from that distance
  \item \emph{Recall:} say \emph{zì — hànzì — hànyǔ} — \textbf{R4}, recognition at distance, eighty lessons back or more
  \item \emph{Recall:} say \emph{zhōngxué — zhōngxuésheng — tóngxué — shàngxué} — \textbf{R4}, recognition at distance, a second word from that far back
\end{itemize}
\end{quote}

\subsection*{Script — a box on the left, a hook on the right}
\begin{quote}
\zh{的}
\end{quote}

\textbf{Eight strokes, in two halves.} The left half is \zh{白} \emph{bái}, \textquotedblleft{}white\textquotedblright{}; the right half is \zh{勺} \emph{sháo}, a ladle. Finish the whole left half before the right one starts — the rule every two-part character in this book has followed.

\textbf{The left half, \zh{白}, five strokes:}

\begin{enumerate}
\raggedright
  \item a short \textbf{piě} falling to the left, at the top
  \item the \textbf{shù} down the left side of the box
  \item the \textbf{héngzhé} across the top and down the right side, without lifting
  \item the \textbf{héng} inside the box
  \item the \textbf{héng} that closes its floor
\end{enumerate}

That box is built exactly as \zh{百}'s was, and for the same reason: \textbf{outside before inside, and the floor last.}

\textbf{The right half, \zh{勺}, three strokes:}

\begin{enumerate}
\raggedright
  \item a \textbf{piě} falling to the left from the top
  \item a long \textbf{héngzhégōu} down the right side, hooking left at the base
  \item a single \textbf{diǎn} dot, placed last, inside the enclosure
\end{enumerate}

\textbf{The dot goes in last.} Enclose, then fill. That is why it is stroke eight and not stroke six.

\subsection*{Your turn}
\begin{itemize}
\raggedright
  \item \emph{Write it:} \zh{的} — \zh{白} first, all five, then \zh{勺}'s three
  \item \emph{Cover:} write it again from memory, and check the dot came last
\end{itemize}

\begin{tcolorbox}[breakable,colback=teal!4,colframe=teal!35!black,title={Before you move on}]
How many strokes does \zh{的} have, and how do they split? (\textbf{Eight} — five for \zh{白}, three for \zh{勺}.) Which stroke is the dot? (\textbf{The last one} — enclose, then fill.) Which character built its box the same way? (\textbf{\zh{百}}.)

Source: \href{https://raw.githubusercontent.com/chanind/hanzi-writer-data/68d10a4b21150cae5e1ebbd223eed289cf32d90c/data/\%E7\%9A\%84.json}{Hanzi Writer Data \zh{的}, pinned stroke paths and medians}.
\end{tcolorbox}
\section[de]{\zh{的} — whose it is}

\label{lesson:ZH-C23-de}



\begin{quote}
One lesson back, and the character is already under your pen.

\begin{itemize}
\raggedright
  \item \emph{Recall:} say \emph{de} — \textbf{R1}, one lesson back, before it decays
  \item \emph{Recall:} say \emph{wǒ shì} — \textbf{R2}, the first real retrieval, five to fifteen lessons back
  \item \emph{Recall:} say \emph{ma} — \textbf{R3}, durable, twenty to sixty lessons back
  \item \emph{Recall:} say \emph{ma} — \textbf{R3}, durable again, a second word from that distance
  \item \emph{Recall:} say \emph{zhōngwén} — \textbf{R4}, recognition at distance, eighty lessons back or more
  \item \emph{Recall:} say \emph{shàngxué} — \textbf{R4}, recognition at distance, a second word from that far back
\end{itemize}
\end{quote}

\begin{sounds}
\textbf{\zh{的}} is said \textbf{de} — and that is the whole of it. \textbf{No tone at all.}

It is the third toneless word in this book, after \zh{吗} and the \zh{上} of \zh{早上}, and like them it is said light, short and low, leaning on the word in front of it. Do not give it a falling tone because it looks like a full character. It is a hinge, not a word.
\end{sounds}

\begin{grammarlens}[title={Grammar Lens: owner + \zh{的} + owned}]
\begin{quote}
\textbf{owner + \zh{的} + the thing owned.}
\end{quote}

\begin{quote}
\textbf{\zh{我的书}} \emph{wǒ de shū} — my book \textbf{\zh{我的名字}} \emph{wǒ de míngzi} — my name \textbf{\zh{你的书}} \emph{nǐ de shū} — your book
\end{quote}

That is the entire Mandarin possessive. Look at what is NOT there:

\begin{itemize}
\raggedright
  \item \textbf{No agreement.} \zh{的} never changes for the owner, the thing, or how many.
  \item \textbf{No separate word for \textquotedblleft{}mine\textquotedblright{}.} English needs \emph{my} and \emph{mine}; Spanish needs \emph{mi}, \emph{mío}, \emph{mía}, \emph{míos}, \emph{mías}. Mandarin has \zh{的}.
  \item \textbf{No difference between one owner and several.} The same \zh{的} does both.
\end{itemize}

You have had \zh{我} since the chapter that taught you to say who you are, and it has had nothing to own until now.

One habit worth learning early: with family and with very close things, Mandarin \textbf{drops} \zh{的} — a speaker says \zh{我家}, not \zh{我的家}. The closer the tie, the likelier the \zh{的} falls away.
\end{grammarlens}

\subsection*{Your turn}
Say these aloud:

\begin{itemize}
\raggedright
  \item \emph{wǒ de shū} — \zh{我的书}
  \item \emph{wǒ de míngzi} — \zh{我的名字}
  \item \emph{nǐ de shū} — \zh{你的书}
\end{itemize}

\begin{tcolorbox}[breakable,colback=teal!4,colframe=teal!35!black,title={Before you move on}]
What tone does \zh{的} carry? (\textbf{None}.) What is the order? (\textbf{Owner, \zh{的}, thing owned}.) And what does \zh{的} do when the owner is plural? (\textbf{Nothing} — it never changes.)
\end{tcolorbox}
\section[kàn shū de rén]{\zh{看书的人} — the one who reads}

\label{lesson:ZH-C23-de-mingci}



\begin{quote}
One lesson back, and the same character doing a bigger job.

\begin{itemize}
\raggedright
  \item \emph{Recall:} say \emph{de} — \textbf{R1}, one lesson back, before it decays
  \item \emph{Recall:} say \emph{shū} — \textbf{R2}, the first real retrieval, five to fifteen lessons back
  \item \emph{Recall:} say \emph{ma} — \textbf{R3}, durable, twenty to sixty lessons back
  \item \emph{Recall:} say \emph{nǐ hǎo ma} — \textbf{R3}, durable again, a second word from that distance
  \item \emph{Recall:} say \emph{wén} — \textbf{R4}, recognition at distance, eighty lessons back or more
  \item \emph{Recall:} say \emph{shàngxué} — \textbf{R4}, recognition at distance, a second word from that far back
\end{itemize}
\end{quote}

\begin{grammarlens}[title={Grammar Lens: put the description in front}]
\zh{的} does not only join an owner to a thing. Put a whole clause in front of it and that clause describes the noun that follows.

\begin{quote}
\textbf{\zh{看书的人}} \emph{kàn shū de rén} — the person \textbf{who reads books} \textbf{\zh{看书的学生}} — the student \textbf{who reads books} \textbf{\zh{我看的书}} \emph{wǒ kàn de shū} — the book \textbf{I read}
\end{quote}

\textbf{Read the Mandarin order against the English.} English says \emph{the person who reads books}: noun first, description trailing after it. Mandarin says \emph{reads-books \zh{的} person}: \textbf{description first, noun last, always.}

That is the single most important thing in this lesson, and it is not a word to learn — it is a direction to turn round. Everything that describes a noun in Mandarin goes in front of it, however long it gets, and \zh{的} is the hinge that holds it there. English has one construction like this and it is short — \emph{a well-read person} — and Mandarin uses that shape for all of them.

Notice that the \zh{的} is the same \zh{的} as the last lesson's, unchanged and toneless. Mandarin does not have one word for the possessive and another for the relative. It has this one hinge, and what goes in front of it decides which job it is doing.
\end{grammarlens}

\subsection*{Your turn}
Say these aloud:

\begin{itemize}
\raggedright
  \item \emph{kàn shū de rén} — \zh{看书的人}
  \item \emph{wǒ kàn de shū} — \zh{我看的书}
  \item \zh{看书的学生}
\end{itemize}

\begin{tcolorbox}[breakable,colback=teal!4,colframe=teal!35!black,title={Before you move on}]
Where does a Mandarin description go? (\textbf{In front of the noun}, always.) Which word holds it there? (\textbf{\zh{的}}.) Is it a different \zh{的} from \zh{我的书}? (\textbf{No} — one hinge, two jobs.)
\end{tcolorbox}
\section[Hànyǔ, wǒ kàn shū]{\zh{汉语} — \zh{我看书}: topic first, comment after}

\label{lesson:ZH-C23-huati}



\begin{quote}
Four in, and this one names the shape the whole language prefers.

\begin{itemize}
\raggedright
  \item \emph{Recall:} say \emph{kàn shū de rén} — \textbf{R1}, one lesson back, before it decays
  \item \emph{Recall:} say \emph{Zhōngguó rén} — \textbf{R2}, the first real retrieval, five to fifteen lessons back
  \item \emph{Recall:} say \emph{nǐ hǎo ma} — \textbf{R3}, durable, twenty to sixty lessons back
  \item \emph{Recall:} say \emph{shì ma} — \textbf{R3}, durable again, a second word from that distance
  \item \emph{Recall:} say \emph{wén} — \textbf{R4}, recognition at distance, eighty lessons back or more
  \item \emph{Recall:} say \emph{tóngxué} — \textbf{R4}, recognition at distance, a second word from that far back
\end{itemize}
\end{quote}

\begin{grammarlens}[title={Grammar Lens: the topic comes first}]
\begin{quote}
\textbf{Mandarin puts the thing being talked about first, and then says something about it.} The first slot is the TOPIC, not necessarily the subject.
\end{quote}

\begin{quote}
\textbf{\zh{汉语} — \zh{我看书}} \emph{Hànyǔ, wǒ kàn shū} — Chinese? I read it. \textbf{\zh{书} — \zh{我看}} \emph{shū, wǒ kàn} — Books, I read. \textbf{\zh{我} — \zh{中国人}} \emph{wǒ, Zhōngguó rén} — Me, I'm Chinese.
\end{quote}

\emph{(Chinese writes a comma at that seam. This book has not taught a single punctuation mark — its own inventory reports the column empty — so the seam is printed here as a dash rather than smuggled in.)}

In each of those the first word is not doing the verb. It is the subject \emph{under discussion}, and the rest of the sentence is the remark about it.

English can do this and treats it as informal — \emph{me, I'd rather walk}; \emph{that book, I've read}. Mandarin does it constantly and treats it as ordinary. It is why so much beginner Mandarin built on an English template sounds stiff: the learner keeps the doer in front when a Mandarin speaker would have put the topic there.

\textbf{A useful consequence.} Once the topic is set, everything afterwards can leave it out. A Mandarin speaker will not repeat \zh{书} in the second half of \emph{shū, wǒ kàn} — and will not supply a pronoun for it either. Where English needs \emph{it}, Mandarin needs nothing, because the topic is still standing.
\end{grammarlens}

\subsection*{Your turn}
Say these aloud:

\begin{itemize}
\raggedright
  \item \zh{汉语} — \zh{我看书}
  \item \zh{书} — \zh{我看}
  \item \zh{我} — \zh{中国人}
\end{itemize}

\begin{tcolorbox}[breakable,colback=teal!4,colframe=teal!35!black,title={Before you move on}]
What goes in the first slot of a Mandarin sentence? (\textbf{The topic} — what is being talked about.) Must it be the one doing the verb? (\textbf{No}.) What does the second half do about repeating it? (\textbf{Nothing} — it leaves it out.)
\end{tcolorbox}
\section[Practice]{\zh{我的书} — \zh{看书的人}}

\label{lesson:ZH-C23-practice}



\begin{quote}
The chapter closes on one character and the four skills.

\begin{itemize}
\raggedright
  \item \emph{Recall:} say \emph{Hànyǔ, wǒ kàn shū} — \textbf{R1}, one lesson back, before it decays
  \item \emph{Recall:} say \emph{Zhōngguó rén — shū — kàn} — \textbf{R2}, the first real retrieval, five to fifteen lessons back
  \item \emph{Recall:} say \emph{shì ma} — \textbf{R3}, durable, twenty to sixty lessons back
  \item \emph{Recall:} say \emph{hǎo bu hǎo} — \textbf{R3}, durable again, a second word from that distance
  \item \emph{Recall:} say \emph{wén} — \textbf{R4}, recognition at distance, eighty lessons back or more
  \item \emph{Recall:} say \emph{tóng} — \textbf{R4}, recognition at distance, a second word from that far back
\end{itemize}
\end{quote}

\begin{grammarlens}[title={What you've built — the hinge}]
One character, eight strokes, no tone — and with it:

\begin{enumerate}
\raggedright
  \item \textbf{\zh{我的书}} — whose something is, with no agreement anywhere.
  \item \textbf{\zh{看书的人}} — a description standing in FRONT of what it describes.
  \item \textbf{\zh{汉语} — \zh{我看书}} — a topic first, and a comment after it.
\end{enumerate}

The third of those needed no character at all: it is an order, not a word. It is in this chapter because \zh{的} is what makes the first slot worth filling.

\textbf{Listen.} \emph{Wǒ de shū. Kàn shū de rén. Hànyǔ, wǒ kàn shū.}

\textbf{Speak.} Say whose the book is, then say who the reader is.

\textbf{Read.} \zh{我的名字} — \zh{看书的学生} — \zh{书} — \zh{我看}

\textbf{Write.} \zh{的} — \zh{白} first, then \zh{勺}, and the dot last.
\end{grammarlens}

\subsection*{Your turn}
\begin{itemize}
\raggedright
  \item \emph{Say it:} \zh{我的书} — \zh{我的名字} — \zh{你的书}
  \item \emph{Say it:} \zh{看书的人} — \zh{看书的学生}
  \item \emph{Write it:} \zh{的}
\end{itemize}

\begin{tcolorbox}[breakable,colback=teal!4,colframe=teal!35!black,title={Before you move on}]
Name the three jobs this chapter gave \zh{的} and its slot. (\textbf{Possession; a description in front of its noun; and the topic that opens a sentence}.) Which one cost no character? (\textbf{The topic} — it is an order, not a word.)
\end{tcolorbox}
